\section{案例簇 A：Robot Soccer 与服务机器人}

前面的应用版图较为抽象，本节用两种 embodied 场景观察同一套 agent stack。Robot Soccer 强调实时多体协作、对抗与角色切换；服务机器人强调感知误差、接触动作和家庭长尾。读图时应关注策略如何落到物理控制，而不只是比赛或演示结果。

\subsection{RoboCup 作为多体实体智能试验场}

RoboCup 把感知、控制和多体策略压缩到一个高速可观测的场景中。这里应观察的不是机器人外形，而是角色如何随球权、队友位置和对手压力实时改变，以及局部动作如何服务于团队目标；这些机制正是后续陌生队友协作的基础。
Stone 用 RoboCup 展示了多 Agent embodied intelligence 的典型难点：感知噪声、动作延迟、策略协同、对手建模、实时决策。机器人需要 ``sensing, deciding, acting'' 的完整闭环，且不能靠远程人工遥控。

\begin{figure}[H]
\centering
\includegraphics[width=0.9\textwidth]{frame-01-robot-soccer.jpg}
\caption{RoboCup 场景：多机器人在真实比赛中自主协同\protect\footnotemark}
\end{figure}

\begin{knowledgebox}{读图：球场是一套动态角色分配系统}
多个机器人同时感知球、队友和对手，并在守门、传球、跑位和抢断之间切换。截图支持“多体策略必须实时协调”，但一场比赛画面不能证明策略能与任意陌生队友协作；这正是后文 ad hoc teamwork 要补上的能力。
\end{knowledgebox}
\footnotetext{视频时间区间：00:08:12--00:10:20。}

\begin{knowledgebox}{为什么 RoboCup 仍有研究价值}
它把多个 AI 子问题耦合在一起：
\begin{itemize}
    \item 局部感知与全局态势理解；
    \item 实时控制与长期战术平衡；
    \item 队内协作与对手对抗；
    \item 规则约束下的高强度策略竞争。
\end{itemize}
\end{knowledgebox}

课程里提到 RoboCup 的长期目标：2050 年 humanoid team 击败人类世界杯冠军队。这个目标是否按时达成并不重要，重要的是它强迫研究者把 ``可发表'' 转化为 ``可比赛''，从而持续暴露系统短板。

\subsection{服务机器人：从竞赛到家庭任务}
另一条线是 RoboCup@Home 场景：做 host、摆台、收纳 groceries、早餐服务等。这些任务难点在于开放环境和长链任务分解，不是单步抓取精度。

\begin{importantbox}{服务机器人任务对 Agent 的要求}
与标准 manipulation benchmark 不同，服务任务需要：
\begin{itemize}
    \item 任务层语义理解（不是只有几何控制）；
    \item 跨步骤记忆与状态追踪；
    \item 出错后的恢复策略；
    \item 与人类交互中的社会可接受行为。
\end{itemize}
\end{importantbox}

\subsection{本章小结}
RoboCup 与服务机器人共同说明：当 Agent 进入实体世界，``策略质量'' 必须同时在感知、控制、协同和规则约束中成立。

